% --- Archon physics formalization source begin ---
% archon:physics
% archon:covers PhyXMiniProblems/problem_phyx_mini_0162.lean
% archon:source-report reports/phyx_mini/problem_phyx_mini_0162.source.json

\chapter{Physics problem phyx\_mini\_0162}
\label{ch:PhyXMiniProblems_problem_phyx_mini_0162}

\paragraph{Problem source.}
\#\# Physical scenario

Figure shows an overhead view of a corridor with a plane mirror \$M\$ mounted at one end. A burglar \$B\$ sneaks along the corridor directly toward the center of the mirror.

\#\# Question

If \$d = 3.0 \textbackslash{}text\{ m\}\$, how far from the mirror will she be when the security guard \$S\$ can first see her in the mirror?

\#\# Answer choices

- A: A: 3.0 m

- B: B: 2.5 m

- C: C: 2.0 m

- D: D: 1.5 m

\#\# Figure caption (auxiliary; use the image as primary evidence)

The image is a diagram depicting a mechanical or structural setup with the following elements:

1. **Structure**: There is an L-shaped structure made of a thicker line, forming a horizontal and vertical section.

2. **Distances**: The diagram shows three distances labeled as "d". Two are horizontal, from the left to the middle section, and from the middle to the end of the horizontal part. The third "d" is vertical along the vertical section of the L-shape.

3. **Points and Labels**:
   - **M**: A point labeled "M" located on the horizontal line left of the middle section.
   - **B**: A point labeled "B" marked below the horizontal line near the left end, with a dashed line connecting it vertically to the structure above.
   - **S**: A point labeled "S" on the vertical section of the L-shape, with a dashed line running vertically from the line to the point.

4. **Shaded Areas**: There are two vertical shaded bars underneath the horizontal structure, not connected to it, with different widths.

The diagram likely represents a physical system or apparatus, potentially dealing with forces, torques, or structural analysis.

\#\# Dataset metadata

Geometrical Optics; Physical Model Grounding Reasoning, Multi-Formula Reasoning

\paragraph{Recorded answer/context.}
D: D: 1.5 m

\paragraph{Figure/image path.}
/root/proposal\_for\_physic/science-mango/phyx\_mini\_run/phyx\_data/test\_image/162.png

\paragraph{Formalization target.}
create a compiling Lean file with sorry bodies at `PhyXMiniProblems/problem\_phyx\_mini\_0162.lean`. The Lean declarations must preserve the physical quantities, dimensions or dimensional roles, figure labels, governing-law hypotheses, and final relation expressed by this problem.
Use Mathlib/Physlib names found through LeanExplore where available. If a physics API is missing, introduce faithful local abstractions rather than scalar placeholder aliases.

\begin{theorem}[Physics formalization target]
\label{thm:physics:phyx_mini_0162:target}
\lean{PhyXMiniProblems.ProblemPhyXMini0162.problem_phyx_mini_0162}
\uses{def:physics:phyx-mini-0162:phyxminiproblems-problemphyxmini0162-lengthquantity, def:physics:phyx-mini-0162:phyxminiproblems-problemphyxmini0162-metres, def:physics:phyx-mini-0162:phyxminiproblems-problemphyxmini0162-physicaldistance, def:physics:phyx-mini-0162:phyxminiproblems-problemphyxmini0162-corridormirrorsetup, def:physics:phyx-mini-0162:phyxminiproblems-problemphyxmini0162-matchescorridorfigure, def:physics:phyx-mini-0162:phyxminiproblems-problemphyxmini0162-obeysfiniteplanemirrorvisibilitylaw, def:physics:phyx-mini-0162:phyxminiproblems-problemphyxmini0162-isfirstvisibledistance, def:physics:phyx-mini-0162:phyxminiproblems-problemphyxmini0162-matchesanswerchoice}
The assigned autoformalize agent should translate this physics problem into Lean declarations in the covered file, with theorem and lemma proof bodies written as `by sorry`.
\end{theorem}
\begin{proof}
This is an autoformalization task, not a proof task. Produce statements that can later be proved by the physics prover without weakening the physical model.
\end{proof}

% --- Archon declaration topology begin ---
\section{Declaration topology}
\begin{definition}[Optical Plane]
  \label{def:physics:phyx-mini-0162:phyxminiproblems-problemphyxmini0162-opticalplane}
  \lean{PhyXMiniProblems.ProblemPhyXMini0162.OpticalPlane}
  Positions in the overhead two-dimensional corridor diagram
\end{definition}
\begin{proof}
  This is the declared model definition of Optical Plane.
\end{proof}
\begin{definition}[Length Quantity]
  \label{def:physics:phyx-mini-0162:phyxminiproblems-problemphyxmini0162-lengthquantity}
  \lean{PhyXMiniProblems.ProblemPhyXMini0162.LengthQuantity}
  A genuine physical quantity with the dimension of length
\end{definition}
\begin{proof}
  This is the declared model definition of Length Quantity.
\end{proof}
\begin{definition}[metres]
  \label{def:physics:phyx-mini-0162:phyxminiproblems-problemphyxmini0162-metres}
  \lean{PhyXMiniProblems.ProblemPhyXMini0162.metres}
  \uses{def:physics:phyx-mini-0162:phyxminiproblems-problemphyxmini0162-lengthquantity}
  Convert an SI-metre readout into a dimensionful physical length
\end{definition}
\begin{proof}
  This is the declared model definition of metres.
\end{proof}
\begin{definition}[length In Metres]
  \label{def:physics:phyx-mini-0162:phyxminiproblems-problemphyxmini0162-lengthinmetres}
  \lean{PhyXMiniProblems.ProblemPhyXMini0162.lengthInMetres}
  \uses{def:physics:phyx-mini-0162:phyxminiproblems-problemphyxmini0162-lengthquantity}
  Read a dimensionful length as a real number of SI metres
\end{definition}
\begin{proof}
  This is the declared model definition of length In Metres.
\end{proof}
\begin{definition}[physical Distance]
  \label{def:physics:phyx-mini-0162:phyxminiproblems-problemphyxmini0162-physicaldistance}
  \lean{PhyXMiniProblems.ProblemPhyXMini0162.physicalDistance}
  \uses{def:physics:phyx-mini-0162:phyxminiproblems-problemphyxmini0162-opticalplane, def:physics:phyx-mini-0162:phyxminiproblems-problemphyxmini0162-lengthquantity, def:physics:phyx-mini-0162:phyxminiproblems-problemphyxmini0162-metres}
  Physical distance between two points whose coordinates are metre readouts
\end{definition}
\begin{proof}
  This is the declared model definition of physical Distance.
\end{proof}
\begin{definition}[Finite Plane Mirror]
  \label{def:physics:phyx-mini-0162:phyxminiproblems-problemphyxmini0162-finiteplanemirror}
  \lean{PhyXMiniProblems.ProblemPhyXMini0162.FinitePlaneMirror}
  \uses{def:physics:phyx-mini-0162:phyxminiproblems-problemphyxmini0162-opticalplane}
  The finite plane mirror marked M in the corridor. The complete affine line carrier supports Mathlib's reflection operation; reflectingSegment is the part on which a physical light ray may strike
\end{definition}
\begin{proof}
  This is the declared model definition of Finite Plane Mirror.
\end{proof}
\begin{definition}[reflect Point]
  \label{def:physics:phyx-mini-0162:phyxminiproblems-problemphyxmini0162-finiteplanemirror-reflectpoint}
  \lean{PhyXMiniProblems.ProblemPhyXMini0162.FinitePlaneMirror.reflectPoint}
  \uses{def:physics:phyx-mini-0162:phyxminiproblems-problemphyxmini0162-opticalplane, def:physics:phyx-mini-0162:phyxminiproblems-problemphyxmini0162-finiteplanemirror}
  The virtual image of a point formed by reflection in the mirror carrier
\end{definition}
\begin{proof}
  This is the declared model definition of reflect Point.
\end{proof}
\begin{definition}[Corridor Mirror Setup]
  \label{def:physics:phyx-mini-0162:phyxminiproblems-problemphyxmini0162-corridormirrorsetup}
  \lean{PhyXMiniProblems.ProblemPhyXMini0162.CorridorMirrorSetup}
  \uses{def:physics:phyx-mini-0162:phyxminiproblems-problemphyxmini0162-opticalplane, def:physics:phyx-mini-0162:phyxminiproblems-problemphyxmini0162-lengthquantity, def:physics:phyx-mini-0162:phyxminiproblems-problemphyxmini0162-finiteplanemirror}
  All physical objects and named landmarks in the overhead corridor model. burglarPositionAt distance is the position labelled B when her distance from the mirror is distance. The remaining four landmarks retain the two horizontal and one vertical spans labelled d in the source description. No field fixes the requested first-visible distance
\end{definition}
\begin{proof}
  This is the declared model definition of Corridor Mirror Setup.
\end{proof}
\begin{definition}[Matches Corridor Figure]
  \label{def:physics:phyx-mini-0162:phyxminiproblems-problemphyxmini0162-matchescorridorfigure}
  \lean{PhyXMiniProblems.ProblemPhyXMini0162.MatchesCorridorFigure}
  \uses{def:physics:phyx-mini-0162:phyxminiproblems-problemphyxmini0162-lengthquantity, def:physics:phyx-mini-0162:phyxminiproblems-problemphyxmini0162-lengthinmetres, def:physics:phyx-mini-0162:phyxminiproblems-problemphyxmini0162-physicaldistance, def:physics:phyx-mini-0162:phyxminiproblems-problemphyxmini0162-corridormirrorsetup}
  The labels and qualitative placement supplied by the corridor description. The burglar moves on a straight line directly toward the mirror center, and the parameter of that path is her physical distance from the center. The three displayed corridor spans are all the dimensionful length d. No visibility threshold or answer value is included here
\end{definition}
\begin{proof}
  This is the declared model definition of Matches Corridor Figure.
\end{proof}
\begin{definition}[Has Unobstructed Image Ray]
  \label{def:physics:phyx-mini-0162:phyxminiproblems-problemphyxmini0162-hasunobstructedimageray}
  \lean{PhyXMiniProblems.ProblemPhyXMini0162.HasUnobstructedImageRay}
  \uses{def:physics:phyx-mini-0162:phyxminiproblems-problemphyxmini0162-opticalplane, def:physics:phyx-mini-0162:phyxminiproblems-problemphyxmini0162-lengthquantity, def:physics:phyx-mini-0162:phyxminiproblems-problemphyxmini0162-finiteplanemirror-reflectpoint, def:physics:phyx-mini-0162:phyxminiproblems-problemphyxmini0162-corridormirrorsetup}
  A method-of-images ray from S to the virtual image of B meets the finite mirror at hit, while both physical ray legs remain in the corridor
\end{definition}
\begin{proof}
  This is the declared model definition of Has Unobstructed Image Ray.
\end{proof}
\begin{definition}[near Horizontal Span In Metres]
  \label{def:physics:phyx-mini-0162:phyxminiproblems-problemphyxmini0162-nearhorizontalspaninmetres}
  \lean{PhyXMiniProblems.ProblemPhyXMini0162.nearHorizontalSpanInMetres}
  \uses{def:physics:phyx-mini-0162:phyxminiproblems-problemphyxmini0162-lengthinmetres, def:physics:phyx-mini-0162:phyxminiproblems-problemphyxmini0162-physicaldistance, def:physics:phyx-mini-0162:phyxminiproblems-problemphyxmini0162-corridormirrorsetup}
  SI-metre readout of the first horizontal corridor span labelled d
\end{definition}
\begin{proof}
  This is the declared model definition of near Horizontal Span In Metres.
\end{proof}
\begin{definition}[far Horizontal Span In Metres]
  \label{def:physics:phyx-mini-0162:phyxminiproblems-problemphyxmini0162-farhorizontalspaninmetres}
  \lean{PhyXMiniProblems.ProblemPhyXMini0162.farHorizontalSpanInMetres}
  \uses{def:physics:phyx-mini-0162:phyxminiproblems-problemphyxmini0162-lengthinmetres, def:physics:phyx-mini-0162:phyxminiproblems-problemphyxmini0162-physicaldistance, def:physics:phyx-mini-0162:phyxminiproblems-problemphyxmini0162-corridormirrorsetup}
  SI-metre readout of the second horizontal corridor span labelled d
\end{definition}
\begin{proof}
  This is the declared model definition of far Horizontal Span In Metres.
\end{proof}
\begin{definition}[vertical Span In Metres]
  \label{def:physics:phyx-mini-0162:phyxminiproblems-problemphyxmini0162-verticalspaninmetres}
  \lean{PhyXMiniProblems.ProblemPhyXMini0162.verticalSpanInMetres}
  \uses{def:physics:phyx-mini-0162:phyxminiproblems-problemphyxmini0162-lengthinmetres, def:physics:phyx-mini-0162:phyxminiproblems-problemphyxmini0162-physicaldistance, def:physics:phyx-mini-0162:phyxminiproblems-problemphyxmini0162-corridormirrorsetup}
  SI-metre readout of the vertical corridor span labelled d
\end{definition}
\begin{proof}
  This is the declared model definition of vertical Span In Metres.
\end{proof}
\begin{definition}[Unfolded Sight Line Clears Wall Edge]
  \label{def:physics:phyx-mini-0162:phyxminiproblems-problemphyxmini0162-unfoldedsightlineclearswalledge}
  \lean{PhyXMiniProblems.ProblemPhyXMini0162.UnfoldedSightLineClearsWallEdge}
  \uses{def:physics:phyx-mini-0162:phyxminiproblems-problemphyxmini0162-lengthquantity, def:physics:phyx-mini-0162:phyxminiproblems-problemphyxmini0162-lengthinmetres, def:physics:phyx-mini-0162:phyxminiproblems-problemphyxmini0162-corridormirrorsetup, def:physics:phyx-mini-0162:phyxminiproblems-problemphyxmini0162-nearhorizontalspaninmetres, def:physics:phyx-mini-0162:phyxminiproblems-problemphyxmini0162-farhorizontalspaninmetres, def:physics:phyx-mini-0162:phyxminiproblems-problemphyxmini0162-verticalspaninmetres}
  After unfolding the reflection in M, the straight sightline clears the occluding wall edge exactly when x * (d\_near + d\_far) ≤ d\_vertical * d\_near. This is the unsolved similar-triangle condition at an arbitrary burglar distance x; it neither selects the first-visible position nor states an answer choice
\end{definition}
\begin{proof}
  This is the declared model definition of Unfolded Sight Line Clears Wall Edge.
\end{proof}
\begin{definition}[Obeys Finite Plane Mirror Visibility Law]
  \label{def:physics:phyx-mini-0162:phyxminiproblems-problemphyxmini0162-obeysfiniteplanemirrorvisibilitylaw}
  \lean{PhyXMiniProblems.ProblemPhyXMini0162.ObeysFinitePlaneMirrorVisibilityLaw}
  \uses{def:physics:phyx-mini-0162:phyxminiproblems-problemphyxmini0162-lengthquantity, def:physics:phyx-mini-0162:phyxminiproblems-problemphyxmini0162-lengthinmetres, def:physics:phyx-mini-0162:phyxminiproblems-problemphyxmini0162-physicaldistance, def:physics:phyx-mini-0162:phyxminiproblems-problemphyxmini0162-finiteplanemirror-reflectpoint, def:physics:phyx-mini-0162:phyxminiproblems-problemphyxmini0162-corridormirrorsetup, def:physics:phyx-mini-0162:phyxminiproblems-problemphyxmini0162-hasunobstructedimageray, def:physics:phyx-mini-0162:phyxminiproblems-problemphyxmini0162-unfoldedsightlineclearswalledge}
  The governing finite-plane-mirror visibility law. It combines specular reflection via the virtual image, obstruction by the corridor walls, and the similar-triangle clearance relation for the captioned L-shaped corridor. Neither Mathlib nor Physlib provides this observer-visibility interface
\end{definition}
\begin{proof}
  This is the declared model definition of Obeys Finite Plane Mirror Visibility Law.
\end{proof}
\begin{definition}[Is First Visible Distance]
  \label{def:physics:phyx-mini-0162:phyxminiproblems-problemphyxmini0162-isfirstvisibledistance}
  \lean{PhyXMiniProblems.ProblemPhyXMini0162.IsFirstVisibleDistance}
  \uses{def:physics:phyx-mini-0162:phyxminiproblems-problemphyxmini0162-lengthquantity, def:physics:phyx-mini-0162:phyxminiproblems-problemphyxmini0162-lengthinmetres, def:physics:phyx-mini-0162:phyxminiproblems-problemphyxmini0162-corridormirrorsetup}
  distance is where S first sees the approaching burglar: she is visible there, but was not visible at any larger nonnegative distance from M
\end{definition}
\begin{proof}
  This is the declared model definition of Is First Visible Distance.
\end{proof}
\begin{definition}[Answer Choice]
  \label{def:physics:phyx-mini-0162:phyxminiproblems-problemphyxmini0162-answerchoice}
  \lean{PhyXMiniProblems.ProblemPhyXMini0162.AnswerChoice}
  The four distance choices printed in the problem
\end{definition}
\begin{proof}
  This is the declared model definition of Answer Choice.
\end{proof}
\begin{definition}[distance In Metres]
  \label{def:physics:phyx-mini-0162:phyxminiproblems-problemphyxmini0162-answerchoice-distanceinmetres}
  \lean{PhyXMiniProblems.ProblemPhyXMini0162.AnswerChoice.distanceInMetres}
  \uses{def:physics:phyx-mini-0162:phyxminiproblems-problemphyxmini0162-answerchoice}
  SI-metre readout printed beside an answer label
\end{definition}
\begin{proof}
  This is the declared model definition of distance In Metres.
\end{proof}
\begin{definition}[Matches Answer Choice]
  \label{def:physics:phyx-mini-0162:phyxminiproblems-problemphyxmini0162-matchesanswerchoice}
  \lean{PhyXMiniProblems.ProblemPhyXMini0162.MatchesAnswerChoice}
  \uses{def:physics:phyx-mini-0162:phyxminiproblems-problemphyxmini0162-lengthquantity, def:physics:phyx-mini-0162:phyxminiproblems-problemphyxmini0162-lengthinmetres, def:physics:phyx-mini-0162:phyxminiproblems-problemphyxmini0162-answerchoice}
  A physical distance agrees with the distance printed for a choice
\end{definition}
\begin{proof}
  This is the declared model definition of Matches Answer Choice.
\end{proof}
\begin{lemma}[first visible distance is half d]
  \label{lem:physics:phyx-mini-0162:phyxminiproblems-problemphyxmini0162-first-visible-distance-is-half-d}
  \lean{PhyXMiniProblems.ProblemPhyXMini0162.first_visible_distance_is_half_d}
  \uses{def:physics:phyx-mini-0162:phyxminiproblems-problemphyxmini0162-lengthquantity, def:physics:phyx-mini-0162:phyxminiproblems-problemphyxmini0162-lengthinmetres, def:physics:phyx-mini-0162:phyxminiproblems-problemphyxmini0162-corridormirrorsetup, def:physics:phyx-mini-0162:phyxminiproblems-problemphyxmini0162-matchescorridorfigure, def:physics:phyx-mini-0162:phyxminiproblems-problemphyxmini0162-obeysfiniteplanemirrorvisibilitylaw, def:physics:phyx-mini-0162:phyxminiproblems-problemphyxmini0162-isfirstvisibledistance}
  The grazing method-of-images ray and the three equal corridor spans give the similar-triangle relation that the first-visible distance is half of d. This is a derived result, not a field of the figure or optics assumptions
\end{lemma}
\begin{proof}
  The conclusion follows from the governing relations and figure readouts named in the statement of first visible distance is half d.
\end{proof}
% --- Archon declaration topology end ---
% --- Archon physics formalization source end ---
